\documentclass[12pt, a4paper]{article}
\usepackage[utf8]{inputenc}
\usepackage[T1]{fontenc}
\usepackage[spanish]{babel}
\usepackage{graphicx}
\usepackage{amsmath}
\usepackage{geometry}
\usepackage{booktabs}
\usepackage{hyperref}
\usepackage{float}
\geometry{margin=2.5cm}

\title{Asimetría del Cambio Climático:\\Inestabilidad Térmica y Enfriamientos Anómalos a Nivel de Países}
\author{
    \textbf{Integrante 1:} Angelo Ibaceta \quad \textbf{Rol:} [Especificar rol] \\
    \textbf{Integrante 2:} [Nombre Integrante 2] \quad \textbf{Rol:} [Especificar rol] \\
    \textbf{Integrante 3:} [Nombre Integrante 3] \quad \textbf{Rol:} [Especificar rol]
}
\date{Abril 2026}

\begin{document}
\maketitle

\section{Sección 1: Análisis de Inestabilidad Térmica (Angelo Ibaceta)}

\subsection{Criterios seleccionados}

Para este trabajo, se eligieron dos criterios cuantitativos que permiten medir la inestabilidad del clima. Esto ayuda a demostrar la idea central del grupo: que el cambio climático no es solo calentamiento, sino también una volatilidad térmica extrema.

\subsubsection{Limpieza y preparación de datos}

Antes de calcular los criterios, el dataset original (\textit{GlobalLandTemperaturesByCountry.csv}) fue sometido a un proceso de análisis y limpieza para garantizar la calidad de los resultados. Se descartaron todos los registros anteriores a 1850 debido a su altísima incertidumbre (propios de la época de termómetros de mercurio tempranos). También se excluyó a la Antártida, ya que el 100\% de sus registros eran nulos. Finalmente, se eliminaron las filas con datos faltantes (\textit{NaN}) en la temperatura y aquellas mediciones con una incertidumbre anómalamente alta (mayor a 3.42~°C, correspondiente a la media más dos desviaciones estándar del período válido). Este filtrado riguroso, reduciendo la base de datos en un 21.8\%, resultó en un conjunto de datos mucho más robusto y confiable.

\subsubsection{Criterio 1: Amplitud térmica anual media}

En primer lugar, se calculó la amplitud térmica para cada país $p$ y año $y$ entre 1990 y 2013. Básicamente, corresponde a la diferencia entre la temperatura promedio mensual máxima y la mínima de cada año:

\[
A_{p,y} = T_{\max}^{(p,y)} - T_{\min}^{(p,y)}
\]

Luego, se obtuvo el promedio histórico de esa amplitud para cada país:

\[
\bar{A}_p = \frac{1}{N_y} \sum_{y} A_{p,y}
\]

Se seleccionó este criterio porque permite detectar qué países sufren los cambios de temperatura más bruscos entre estaciones. Un valor alto en este indicador señala que el régimen térmico de la zona se está volviendo más inestable y difícil de predecir.

\subsubsection{Criterio 2: Anomalías severas de enfriamiento}

El segundo criterio busca detectar extremos de frío que están fuera de los rangos normales. Se tomó el período de 1950 a 1989 como línea base para calcular el promedio $\mu$ y la desviación estándar $\sigma$. Luego, se usaron los datos recientes (1990--2013) para calcular el Z-score:

\[
Z_{p,m,y} = \frac{T_{p,m,y} - \mu_{p,m}}{\sigma_{p,m}}
\]

Se definió como ``enfriamiento anómalo severo'' cualquier mes donde el Z-score fuera menor a $-2$. La severidad total por país ($S_p$) se midió como el promedio de estos eventos extremos:

\[
S_p = \frac{1}{|C_p|} \sum_{(m,y) \in C_p} |Z_{p,m,y}|
\]

Este punto es fundamental para el análisis grupal: demuestra que aunque el planeta se caliente en términos globales, algunas regiones sufren golpes de frío estadísticamente extremos, lo que prueba que el sistema climático es cada vez más errático.

\subsubsection{Intersección de criterios}

Para seleccionar los 8 países finales, se normalizaron los resultados de ambos criterios y se sumaron los puntajes dándoles el mismo peso (50\% a cada uno):

\[
\text{Score}_p = 0.5 \cdot \hat{A}_p + 0.5 \cdot \hat{S}_p
\]

De esta forma, el gráfico resultante muestra a los países que presentan simultáneamente las amplitudes térmicas más altas y los episodios de frío más anómalos.

\subsection{Gráfico elaborado y fuente de datos}

\begin{figure}[H]
    \centering
    \includegraphics[width=\textwidth]{grafico_asimetria.png}
    \caption{Gráfico tipo Dumbbell que ilustra la relación entre la amplitud térmica extrema y la severidad de episodios de enfriamiento (1990--2013). Las temperaturas promedio mínimas y máximas se representan mediante círculos y rombos, respectivamente. La intensidad del color denota la magnitud de las anomalías de frío, mientras que la etiqueta $n$ contabiliza la frecuencia de dichos eventos.}
    \label{fig:dumbbell}
\end{figure}

\textbf{Fuente de datos:} Climate Change: Earth Surface Temperature Data de Berkeley Earth, extraído de Kaggle (\url{https://www.kaggle.com/datasets/berkeleyearth/climate-change-earth-surface-temperature-data}).

\subsection{Conclusiones obtenidas}

A partir de la visualización generada, es posible extraer observaciones clave sobre la inestabilidad climática en la región euroasiática. En primer lugar, la longitud de los segmentos en el gráfico evidencia la magnitud de la amplitud térmica de cada país. Rusia y Mongolia destacan notoriamente con las barras más largas, indicando variaciones anuales extremas que promedian más de 40~°C de diferencia entre sus temperaturas máximas y mínimas. 

Sin embargo, el uso de la escala de color revela un fenómeno aún más complejo, los países con mayor amplitud térmica no son necesariamente los que experimentan los enfriamientos más severos. Al observar los tonos más intensos en el gráfico, resalta el caso de Tayikistán, que presenta la anomalía de frío más extrema ($S_p = 2.99$) pese a tener un segmento visiblemente más corto que el de Rusia (27.0~°C de amplitud). Del mismo modo, países como Estonia y Bielorrusia, aunque tienen variaciones térmicas más moderadas, concentran una mayor cantidad de olas de frío, lo cual se refleja en las etiquetas ($n$) que acompañan a sus barras.

En síntesis, esta representación visual respalda de manera clara la hipótesis central del estudio. El cambio climático no se manifiesta únicamente a través de un calentamiento constante, sino que se traduce en una profunda desestabilización del sistema. Las enormes oscilaciones térmicas y los episodios de frío atípico, ilustrados mediante la longitud y el color en el gráfico, demuestran que la volatilidad es un indicador crítico de la actual crisis climática.

\newpage
\section{Sección 2: [Título del Análisis] ([Nombre Integrante 2])}
% -----------------------------------------------------------------------------
% Instrucciones Integrante 2:
% Reemplaza el título y tu nombre en el comando \section de arriba.
% Redacta todo tu contenido, análisis, y agrega tus subsecciones o gráficos aquí.
% -----------------------------------------------------------------------------

[El integrante 2 debe redactar su parte aquí...]

\newpage
\section{Sección 3: [Título del Análisis] ([Nombre Integrante 3])}
% -----------------------------------------------------------------------------
% Instrucciones Integrante 3:
% Reemplaza el título y tu nombre en el comando \section de arriba.
% Redacta todo tu contenido, análisis, y agrega tus subsecciones o gráficos aquí.
% -----------------------------------------------------------------------------

[El integrante 3 debe redactar su parte aquí...]

\newpage
\section{Sección 4: Visualización generada por Inteligencia Artificial}
% -----------------------------------------------------------------------------
% Instrucciones Generales:
% En esta sección final se debe incrustar la visualización que haya sido 
% generada íntegramente por IA, junto con su respectiva descripción o análisis.
% -----------------------------------------------------------------------------

[Incluir aquí la visualización de IA y el análisis correspondiente...]

\end{document}